\begin{explanationbox}
\textbf{一、先把椭圆看成一个被横向拉长的圆。}

椭圆上每个点 $P(x,y)$ 都满足 $x^2/a^2+y^2/b^2=1$。如果换坐标
\[
 u=\frac{x}{a},\qquad v=\frac{y}{b},
\]
椭圆就变成单位圆 $u^2+v^2=1$。圆在 $U(u_0,v_0)$ 处的切线与半径 $OU$ 垂直，因此
\[
 u_0(u-u_0)+v_0(v-v_0)=0
 \quad\Longrightarrow\quad u_0u+v_0v=1.
\]
把 $u=x/a$、$u_0=x_0/a$ 等代回去，立刻得到椭圆切线式。\textbf{理解关键：并非凭空把平方“拆开”，而是圆的切线经过缩放后的表达。}

\medskip
\textbf{二、为什么切线只碰一下，不穿过椭圆？}

设 $P(x_0,y_0)$ 是切点，任取椭圆上一点 $Q(x,y)$。由柯西不等式可得
\[
 \frac{x_0x}{a^2}+\frac{y_0y}{b^2}
 \le\sqrt{\frac{x_0^2}{a^2}+\frac{y_0^2}{b^2}}
     \sqrt{\frac{x^2}{a^2}+\frac{y^2}{b^2}}=1.
\]
等号当且仅当 $Q=P$。这表明整条椭圆位于直线
$x_0x/a^2+y_0y/b^2=1$ 的同一侧，只在 $P$ 处接触。\textbf{切线是贴着椭圆的一条支撑线，而不是穿过椭圆的割线。}

\medskip
\textbf{三、为什么需要法向量？斜率不够用吗？}

直线 $Ax+By=C$ 的法向量是 $(A,B)$，它\textbf{垂直于直线}。所以切点式一写出来，便能直接读出
\[
 \boldsymbol n=\left(\frac{x_0}{a^2},\frac{y_0}{b^2}\right).
\]
它能统一处理\textbf{竖直切线、斜率、平行和垂直}，不用先考虑 $y_0$ 是否为零。法向量一般\textbf{不是}从原点指向切点的向量 $(x_0,y_0)$：椭圆横纵方向的缩放不同。只有在特殊位置（如长短轴端点）或圆的情形，两者才自然同向。

\medskip
\textbf{四、外点 $M$ 为何能生成一条“切点连线”？}

先想象从椭圆外点 $M$ 拉出两条绷紧的线，刚好分别碰到椭圆，接触位置是 $T_1$、$T_2$。它们形成一条\textbf{联系两个切点的直线}，而不是经过 $M$ 的第三条切线。

设其中一个切点为 $T(u,v)$。因为 $M(x_1,y_1)$ 落在 $T$ 处的切线上，所以
\[
 \frac{ux_1}{a^2}+\frac{vy_1}{b^2}=1.
\]
把 $(u,v)$ 看成待求的\textbf{切点坐标} $(x,y)$，就是
\[
 \frac{x_1x}{a^2}+\frac{y_1y}{b^2}=1.
\]
两个切点都满足这条直线方程，故它们的连线就是极线。\textbf{同一个代换公式，已知的是椭圆上的点时求切线；已知的是外点时求切点弦。}

\medskip
\textbf{五、用“缩放后的圆”理解外点的三种位置。}

把 $M$ 也换成 $U=(x_1/a,y_1/b)$，则 $OU^2=S$。极线在圆坐标中变为
\[
 \frac{x_1}{a}u+\frac{y_1}{b}v=1.
\]
当 $S>1$，它到圆心的距离为 $1/\sqrt S<1$，因此与圆有两个交点，即两个实切点；当 $S=1$，只有一个交点；当 $0<S<1$，距离大于 $1$，不相交。$M=O$ 时变为 $0=1$，不能当成一般直线。

\medskip
\textbf{六、考场上如何最快记住、最快用？}
\begin{itemize}
    \item \textbf{口诀：平方拆成“切点乘变量”，分母原封不动。} 例如 $x^2/a^2$ 换成 $x_0x/a^2$。
    \item \textbf{先看点的位置：} 点在椭圆上，写切线；点在椭圆外，写切点弦；点在椭圆内，没有两条实切线。
    \item \textbf{先写直线，后求斜率：} 统一使用切点式，特别能避免 $y_0=0$ 时的除零错误。
    \item \textbf{辨认试题信号：} “已知切点”“过外点两切线”“切点连线”“斜率”“相切判定”，优先想到本结论。
\end{itemize}
\end{explanationbox}
